\chapter{Introduction}\label{ch:intro}

\section{Broad Area of Work}\label{sec:area}
The broad area of this dissertation is Applied Natural Language Processing
(NLP) and multimodal machine learning for the analysis of public retail
customer feedback on social media.  The work spans four technical sub-areas:
(a) supervised text classification with transformer encoders for sentiment
in short, informal English social-media text;
(b) zero-shot natural-language-inference (NLI) models for aspect-based opinion
mining over a fixed retail taxonomy;
(c) multimodal vision-language models for captioning and OCR of screenshots
and photographs attached to retail complaints; and
(d) lightweight trust and credibility filtering using account-metadata
heuristics, near-duplicate detection, and retail-insider terminology signals.
Around these models sits a data-engineering layer for scheduled ingestion of
public posts and structured storage suitable for aggregation, and a serving
layer that surfaces the results in an interactive dashboard.

\section{Motivation}\label{sec:motiv}
Customers, associates, and delivery drivers post continuously about their
Walmart, Sam's Club, Spark-driver, and OGP/pickup experiences on Reddit.  In
the author's role at Walmart Global Tech this signal is presently consumed
through ad-hoc browsing and keyword-filter dashboards, which suffer from poor
coverage, high noise, and latency: a complaint that trends today may not
surface in a manual review for days.  Conventional lexicon- or rule-based
sentiment systems perform poorly on the sarcasm, slang, and long-form
narrative style typical of Reddit retail posts, and none of them filter fake
or low-credibility content before aggregation.

This dissertation demonstrates that a stack of task-specific models --- a
fine-tuned encoder for sentiment, a zero-shot NLI model for aspects, and a
multimodal vision model for images --- combined with an explicit trust filter
can convert a raw public stream into a structured, aspect-tagged,
trust-weighted feed suitable for operational review.

\section{Contributions of this Dissertation}\label{sec:contrib}
The principal contributions of the completed work are:
\begin{enumerate}
\item A modular, offline-first, zero-API-cost end-to-end pipeline (Retail
Sentiment Intelligence) covering ingestion, trust scoring, sentiment,
aspects, vision, aggregation, alerting, and dashboarding.
\item A domain fine-tuned \textbf{ModernBERT} sentiment classifier trained
through a three-stage curriculum, raising macro-F1 from 0.6272 to
\textbf{0.7642} overall and from 0.2778 to \textbf{1.0000} on long posts,
evaluated with 5-fold out-of-fold cross-validation.
\item A multi-pass vision-captioning technique that adapts ideas from recent
vision-language papers onto a compliant Gemma~3 4B model, reducing
hallucination from 50\% to \textbf{0\%} and lifting text extraction from 25\%
to \textbf{75\%}.
\item An interpretable trust score and a \textsf{trust $\times$ confidence}
admission gate that flags --- rather than drops --- low-credibility posts,
with every constant mapped to a stakeholder-arguable English rationale.
\item A React/FastAPI dashboard (Brand Health, Alert Feed, Post Explorer,
Review \& Validate, Post Lifecycle, Insights \& Competitor, Pipeline
Control, Notifications) with a human-in-the-loop Review \& Validate workflow
whose corrections and posted replies feed few-shot reply generation and a
roadmap to automatic replies.
\end{enumerate}

\section{Organisation of the Report}\label{sec:org}
Chapter~\ref{ch:problem} states the problem and objectives.
Chapter~\ref{ch:lit} surveys the relevant literature.
Chapter~\ref{ch:design} presents the layered system design and end-to-end
pipeline.  Chapter~\ref{ch:impl} describes the implementation of each module.
Chapter~\ref{ch:eval} reports the experimental evaluation and results.
Chapter~\ref{ch:tools} documents the tools and configuration used.
Chapter~\ref{ch:problems} records the problems encountered and mitigations.
Chapter~\ref{ch:concl} draws conclusions and Chapter~\ref{ch:future} outlines
future work.  A glossary, references, and supporting appendices close the
report.
